% answers_short.tex —— 2018 年·北京卷·理科 · 简短答案速查
% 由 tools/make_short_answers.py 自 answers.tex 的 \gksolhead/\gkanswer 提取，
% 并与 2018/2018_数学真题_北京卷理_答案要点.md 逐题比对后生成，勿手改（改源 MD 后重跑）。

\begin{gknotice}
\textbf{使用说明}：本节只给最终答案便于快速核对，完整推导见解析版。答案为编写者独立推导，\textbf{非官方标准答案}。
\end{gknotice}

\gkpart{选择题}
\begin{center}\begin{tabular}{ccccccccc}
\toprule
题号 & 1 & 2 & 3 & 4 & 5 & 6 & 7 & 8 \\
\midrule
答案 & A & D & B & D & C & C & C & D \\
\bottomrule
\end{tabular}\end{center}

\gkpart{填空题}
\begin{itemize}[leftmargin=2.6em]
  \item[\textbf{9.}] $a_n=6n-3$
  \item[\textbf{10.}] $1+\sqrt2$
  \item[\textbf{11.}] $2/3$
  \item[\textbf{12.}] $3$
  \item[\textbf{13.}] 例如 $f(x)=x(3-x)$
  \item[\textbf{14.}] 椭圆 $\sqrt3-1$；双曲线 $2$
\end{itemize}

\gkpart{解答题}
\begin{itemize}[leftmargin=2.6em]
  \item[\textbf{15.}] (1) $\angle A=\pi/3$；(2) $3\sqrt3/2$
  \item[\textbf{16.}] (1) 待核；若假设直三棱柱且侧棱长2，则条件性结论成立；(2) 待核；若假设直三棱柱且侧棱长2，则余弦为 $-1/\sqrt{21}$；(3) 待核；若假设直三棱柱且侧棱长2，则条件性结论成立
  \item[\textbf{17.}] (1) $1/40$；(2) $7/20$；(3) $D\xi_1>D\xi_4>D\xi_2=D\xi_5>D\xi_3>D\xi_6$
  \item[\textbf{18.}] (1) $a=1$；(2) $a>1/2$
  \item[\textbf{19.}] (1) $k<1$ 且 $k\ne0,-3$；(2) $1/\lambda+1/\mu=2$
  \item[\textbf{20.}] (1) $2,1$；(2) 最大 $4$；(3) 例如 $\{0,e_1,\ldots,e_n\}$，最大 $n+1$
\end{itemize}
